%% ============================================================
%% Master's Thesis Defense — Speaker notes
%% Nidanur Günay — what to say, slide by slide
%% Compile from the PROJECT ROOT:
%%   pdflatex -interaction=nonstopmode Defense_Speaker_Notes.tex
%% ============================================================

\documentclass[11pt,a4paper]{article}

\usepackage[T1]{fontenc}
\usepackage[utf8]{inputenc}
\usepackage[english]{babel}
\usepackage[a4paper,top=2.0cm,bottom=2.0cm,left=2.2cm,right=2.2cm]{geometry}
\usepackage{helvet}
\renewcommand{\familydefault}{\sfdefault}
\usepackage{xcolor}
\usepackage{graphicx}
\usepackage{booktabs}
\usepackage{enumitem}
\usepackage{setspace}
\usepackage{needspace}
\usepackage{fancyhdr}
\usepackage[hidelinks]{hyperref}

\definecolor{knblue}{HTML}{00539F}
\definecolor{knlight}{HTML}{EAF1F8}
\definecolor{kngrey}{HTML}{5A6570}
\definecolor{knaccent}{HTML}{B45309}

\setstretch{1.12}
\setlength{\parindent}{0pt}
\setlength{\parskip}{0.55em}

\pagestyle{fancy}
\fancyhf{}
\lhead{\footnotesize\color{kngrey}Defense speaker notes}
\rhead{\footnotesize\color{kngrey}Nidanur Günay}
\cfoot{\footnotesize\color{kngrey}\thepage}
\renewcommand{\headrulewidth}{0.4pt}

%% Slide header: number, title, time
\newcommand{\slidehead}[3]{%
  \needspace{5\baselineskip}
  \vspace{0.9em}
  \noindent\colorbox{knblue}{\parbox{0.978\textwidth}{\vspace{1pt}\color{white}\textbf{Slide #1 \quad #2}\hfill\textbf{#3}\vspace{1pt}}}
  \vspace{0.35em}\par}

%% "Say this" body is just normal text.
%% Reminder box
\newcommand{\remind}[1]{%
  \vspace{0.15em}
  \noindent\colorbox{knlight}{\parbox{0.972\textwidth}{\small\color{knblue}\textbf{Remember:} #1}}
  \vspace{0.2em}\par}

\newcommand{\pointto}[1]{{\color{knaccent}\textit{[#1]}}}

\begin{document}

%% ---------------- Cover ----------------
\begin{center}
  {\LARGE\bfseries\color{knblue} Defense — what to say}\\[0.4em]
  {\large Abstraction Mechanisms for Virtual Avatars}\\[0.3em]
  {\color{kngrey} Nidanur Günay \quad|\quad 20 minutes \quad|\quad 33 slides plus backup}
\end{center}

\vspace{0.6em}

\noindent\colorbox{knlight}{\parbox{0.972\textwidth}{\small\color{knblue}
\textbf{How to use this.} The text under each slide is written the way you would actually say it out loud. It is not meant to be read word for word, it is meant to keep you on track. Every slide has a time budget. If you are running late, the slides marked \textbf{CUT} can be shortened to one sentence without breaking the story.\\[0.4em]
\textbf{Three habits that matter more than the words.} Pause after every result. Say the verdict out loud (\emph{supported}, \emph{not supported}, \emph{inconclusive}) before you explain it. Never apologise for a null result.}}

\vspace{0.9em}

\begin{center}
\begin{tabular}{@{}llr@{}}
\toprule
\textbf{Part} & \textbf{Slides} & \textbf{Time} \\
\midrule
Opening: title, clip, problem & 1--3 & 1:01 \\
\textbf{Related work} & \textbf{4--8} & 3:13 \\
RQs and the system & 9--12 & 1:53 \\
Applying the techniques, texture & 13--14 & 0:55 \\
Pre-study and composite & 15--16 & 1:05 \\
Answering RQ1 and RQ2 & 17--18 & 1:15 \\
Main study, hypotheses, subscales & 19--20 & 1:40 \\
\textbf{Results, H1 to H4} & \textbf{21--24} & 2:40 \\
Behaviour and the RQ3 answer & 25--26 & 1:15 \\
\textbf{Discussion} & \textbf{27--29} & 1:38 \\
Limitations & 30 & 0:35 \\
Conclusion, future work, thanks & 31--33 & 1:45 \\
\midrule
& & \textbf{$\approx$ 18:55} \\
\bottomrule
\end{tabular}
\end{center}

\vspace{0.5em}
\noindent\colorbox{knlight}{\parbox{0.972\textwidth}{\small\color{knblue}
\textbf{Related work is four slides, 5 to 8, and 3:15.} The two overview slides that previously stood in front of them presented the same evidence at a coarser resolution, and they now sit in the backup section with their own scripts at the end of this document. Each of the four that remain carries one strand: the valley, then trust, then what trust actually means, then abstraction.\\[0.4em]
\textbf{If you are running behind}, these three are still the cheapest place to buy time. Slide 5 compresses to ``Mori drew the curve, MacDorman and Ishiguro measured it, Ho and MacDorman built the scale I use, and Dill showed it reaches CG characters.'' Slides 6 and 8 each compress to their blue bar. Leave slide 7 alone, the trust definition is load bearing. That buys about a minute and a half without breaking the argument.}}

\clearpage

%% ============================================================
\section*{\color{knblue}Part 1 — Setup (slides 1 to 12, about 6 minutes)}

\slidehead{1}{Title}{0:10}

Good afternoon everyone, and thank you for being here. My name is Nidanur Günay and today I will present my master's thesis, \emph{Abstraction Mechanisms for Virtual Avatars}.

\remind{That is the whole slide. Do not add what the thesis is about, slide 2 does that. Click straight through.}

\slidehead{2}{Motivation}{0:16  (13 s clip, autostart)}

The clip starts by itself as the slide opens. Let it run about two seconds, then speak over it. Do not wait for it to finish and do not narrate what is on screen.

Virtual avatars are already part of our daily lives, through VR and generative AI. They present, they teach, they advise, and they sell. And how they are drawn decides how much we believe what they say.

\remind{That is about 14 seconds of speech against a 13-second clip, meaning your last few words land just as it ends. The clip is doing the work here. Resist adding examples. Then move straight on.\\[0.4em] The clip is cut from three sources and they are named in small type under it, with the full URLs on the Media credits slide (35). Do not read them out. The attribution has to be visible on screen, not spoken.}

\slidehead{3}{Problem Statement}{0:35}

In real applications you would expect that the more human an avatar looks, the more comfortable we feel with it. That is true, but only up to a point. When a character gets very close to human but not quite there, people start to find it creepy instead.

That dip is the uncanny valley, a well-known phenomenon in CG, proposed by Mori initially for humanoid robots. \pointto{point at the curve} However, later studies revealed that it holds for CG characters as well.

\remind{Do not say \emph{trust} on this slide, and do not mention NPR yet. This slide is only the valley: likeness against liking, which is what Mori's curve actually plots. NPR and trust both arrive on the next slide, where the evidence for them does.}

\slidehead{4}{Related Work}{0:03}

\remind{Section divider. Click straight through, do not read it out.}

\slidehead{5}{1. Measuring the valley}{0:50 \quad CUT}

Mori drew that curve in 1970, but he did not measure it.

MacDorman and Ishiguro did it in 2006. They morphed a robot into a real face step by step, and eeriness peaked exactly where comfort was lowest. Eeriness here is not dislike. It is the unsettled response to something that approaches human likeness without reaching it. The curve now had a second measure, running against comfort.

Ho and MacDorman went further. The scales available at the time could not separate the two axes, so they built one where eeriness is measured independently of human likeness. Without that separation, a drop in eeriness could just mean the character looks less human.

Dill then extended the scope. The same curve appears for CG film and game characters, not only robots, and the eyes and the mouth drive the response.

By 2012, then, the valley is real, it is measurable, and it applies to avatars.

\remind{Two things to hold back. Do not say that you use Ho and MacDorman's instrument, that belongs in the methodology. And do not say \emph{trust} yet, this slide is still comfort and eeriness. No participant counts either.}

\slidehead{6}{2. Trust follows the same curve}{0:45 \quad CUT}

However, the valley reaches further than comfort. It reaches trust.

Nowak and Rauh showed that first on avatars. Avatars judged more human-like were also rated more credible, and they were chosen more often.

Later, McDonnell put trustworthiness on a rating scale and ran it across eleven render styles on two motion-captured actors. The result, though, is valley-shaped, and the single most trustworthy style of all was ToonPencil, the most abstract one they had.

Natarajan and Gombolay found the same pull outside avatars entirely, in human-robot interaction. More anthropomorphic agents were trusted more.

The valley therefore does reach trust. Which leaves one question: how do you measure it?

\remind{Two things. The \emph{though} in the McDonnell sentence is load-bearing, because without it the valley-shaped result sounds like the next item on a list rather than a correction to Nowak and Rauh. Keep it, or say ``but'' instead.\\[0.4em] You have dropped the spoken justification for why trust matters, which is fine, it is on the slide in the blue bar. Let the room read it while you talk.\\[0.4em] End on the question and stop. Do not answer it, the next slide is the answer.}

\slidehead{7}{Measuring Trust}{0:35}

Mayer and colleagues answered that in 1995. They separated two things that often get confused. Trustworthiness is a property of the person being trusted. Trust is the attitude you hold towards them.

And they defined trustworthiness as three dimensions.

Ability. Competence inside one domain, not in general.

Benevolence. The belief that the other party cares about my wellbeing.

And integrity. Whether they stick to principles I find acceptable, and whether what they say matches what they do.

Those three move independently. Someone can be entirely competent and still not act in your interest, which is why one trust score would hide which part actually moved.

\remind{\textbf{``Mayer and colleagues'', then ``they''.} The paper is Mayer, Davis and Schoorman. The plural is simply correct, and it saves you having to guess at anyone's pronouns mid-sentence.\\[0.4em] \textbf{Integrity is not transparency.} Someone can be completely open about having no principles. Mayer's integrity is adherence to principles the \emph{trustor} finds acceptable. ``Sticks to principles I accept'' is therefore the safe phrasing.\\[0.4em] Point at each box in the diagram as you name it.}

\slidehead{8}{3. Is abstraction the exit?}{1:00 \quad CUT}

That model of trust was used again by Canales. They compared three stylisation levels of an avatar doctor and measured trust on exactly those three dimensions, with a forced choice on top. Self-report was flat across all three styles, but when people had to actually pick one, forty-eight per cent chose the semi-realistic doctor. Their strongest effect, though, was not style at all. It was identity, which is why I hold the speaker constant in my own design.

That brought the question this thesis starts from. Can abstraction be a way out of the valley?

Dubosc tested it in an interactive VR study, and found the stylised avatars were rated significantly less eerie.

However, there is no consensus. Gao set a stylised agent against video of the real person it was modelled on, and the real speaker scored higher on all three dimensions. Tao pooled thirteen studies and found high realism rated more trustworthy than low. And Alipour reviewed fifty-three studies and concluded the link between eeriness and trust is simply inconsistent.

The question is open. And nobody has tested one avatar under both treatments against the real person it came from.

\remind{``That model of trust was used again by Canales'' carries the previous slide straight into this one. Do not re-introduce Mayer.\\[0.4em] The identity sentence is the first thing to drop if you are running late. It explains the third line of panel 01, but it also previews a design choice you defend later. Losing it here costs nothing.\\[0.4em] Pause after ``no consensus''. Two studies point one way, two the other, and the review says there is no consistent answer. That pause is what the next slide grows out of.}

\slidehead{9}{Research Questions}{0:40}

That gap is where this thesis starts, and it splits into three questions.

The first is technical. How do these techniques adapt to avatar platforms with completely different rendering architectures, and what capabilities and limitations does each platform impose.

The second is about selection. Which technique, or which combination, best balances visual abstraction against a face and an expression that stay readable.

And the third is perception. How does stylisation affect how eerie and how trustworthy an avatar feels, and can it reach the trustworthiness of a photorealistic avatar, or of a real human speaker.

The first two are answered by the system I built. The third by a user study, and that is where most of this talk goes.

\slidehead{10}{Methodology \& Evaluation}{0:03}

\remind{Section divider. Click straight through, do not read it out.}

\slidehead{11}{Three Avatar Platforms}{0:25}

In order to design the system and answer RQ1, three avatar platforms were chosen. They differ in how much of the renderer they expose and in whether they can animate a face, and each has its own strengths and limits.

Among these platforms, Avaturn is the only one that has a realistic face and can animate that face. Our study avatar was made from a photo of me, as you can see. The user study needed both of those, and Avaturn was therefore the platform selected for it.

However, all three platforms were used to develop and implement the eight NPR techniques, and those are on the next slide.

\remind{Do not walk the table. They can read it, and anything they want to know is a question afterwards.\\[0.4em] ``A photo of me'' is the one light moment in the setup. Let it land, then move on.\\[0.4em] If asked: Mixamo is fully open but its exports have no facial rigging. The Meta SDK renders internally, one hook after its own lighting, and gives back live body tracking in VR. Shader access matters too, but Mixamo has that as well, so it is not what decided the choice.}

\slidehead{12}{Eight Techniques}{0:45}

These eight techniques were developed iteratively, each one resolving a limitation of the technique before it. X-Toon adds control over toon shading. The inverted hull and the normal based edge detection were missing the lines that come from the texture, and that is what the Sobel family answers.

After the development, each technique was assessed in the Unity editor and on the VR headset against four criteria: can you still read the face, is it stable in motion, are the artefacts suppressed, and does it read as a deliberate style rather than a broken render.

In the result of that first assessment, five of the eight were kept as candidates.

\remind{Name the four criteria out loud. This table is an evaluation result and they have not seen any method yet, so ``assessed in the editor and on the headset against four criteria'' is what makes it credible in four seconds.\\[0.4em] Do not say what the composite ends up being. That is slide 18.}

\slidehead{13}{Applying the Techniques}{0:30}

During that initial analysis, the texture based edge detection methods were failing to extract the feature lines on the face.

\pointto{Top figure.} Gaussian prefiltered Sobel. \pointto{Bottom figure.} The hierarchical pass, reading depth, normal and colour together. Both produce noisy output, and neither finds the eyebrows or the lip line.

The reason is the baked light already in the texture of the avatar, which comes from the way that avatar is built. An edge detector cannot tell that lighting from a real feature.

In order to solve this problem, a texture modification was applied.

\remind{Stop there and change slide. The fix is the next slide, do not give it away here.\\[0.4em] If asked why raising the threshold does not help: it removes the noise and the features together. Each panel on screen is a different threshold.\\[0.4em] If asked about the other techniques: normal edges find almost nothing on smooth skin and their line weight grows with camera distance, and Kuwahara abstracts the surface without drawing lines at all. Backup slides cover both.}

\slidehead{14}{Avatar Texture Manipulation}{0:25}

Therefore, to get clearer results and see the real capability of the techniques, we modified the face texture slightly in Blender. Blurring, unifying the skin tone, and raising the saturation a little across the mouth.

Every NPR technique from here on is applied to this modified avatar. However, the photorealistic one keeps its original texture.

\remind{Do not say C2 or C3, the conditions arrive on slide 19. Say \emph{the stylised avatar} and \emph{the photorealistic one}.\\[0.4em] If asked whether it worked: Sobel then finds the lip line and the eyebrows, and the Gaussian prefilter became unnecessary, which is why it is not in the final composite.\\[0.4em] If challenged on this as a confound: extending the preprocessing to the photorealistic condition would destroy the very thing that condition exists to represent. A design decision, not an oversight.}

\slidehead{15}{Parameter Adjustment Study}{0:30}

In order to decide which of the candidate techniques to use, we conducted a parameter adjustment study with 45 participants.

Each round took the parameters of one candidate technique and asked people to rank the renders by how trustworthy the advisor looked.

And in the last question I asked for the overall best. Across everything you have seen, which avatar looked most appealing.

\pointto{The lower chart.} In the results, R1\,E came first with seven votes, R3\,B second with six.

\remind{Do not give any outcome yet. This slide is only what participants did. The results are the next slide.\\[0.4em] \textbf{If you have time in hand}, the two things worth adding back are: each round's winner was frozen for every round after it, meaning the composite was built up one decision at a time; and the strip on screen is one round as an example, the inner line threshold from 0.01 on the left to 0.05 on the right, where the shirt goes from almost solid black to clean lines. Both are already on the slide, so only say them if the clock allows.}

\slidehead{16}{Parameter Adjustment Study}{0:35}

Here you see the manipulated parameters and the results of the study. All three tests are non-parametric. Borda uses the whole ballot, Friedman asks whether a round was ranked at random, and Wilcoxon compares the top two.

\pointto{The table.} The first visible result is that none of the post-processing effects was selected. Neither the Kuwahara filter nor the hierarchical layer survived.

And in order to decide the rest, which had non-significant results in their rounds, both the Borda winner and the last overall appeal question were taken into account.

\remind{Three of the seven final values came from the appeal question rather than a significant ranking. The last column records it, and someone will read it.\\[0.4em] If pushed on the tests: Dem\v{s}ar recommends Friedman for more than two conditions and Wilcoxon for a single pair. You followed a published procedure, you did not pick one.}

%% ============================================================
\section*{\color{knblue}Part 3 — The user study (slides 19 to 20, about 3 and a half minutes)}

\slidehead{17}{RQ1 Answer}{0:35}

After setting up the system, RQ1 and RQ2 were answered as well. The first question asked how these techniques adapt to platforms with different rendering architectures, and what each platform imposes.

All eight techniques ran on all three platforms, with a different implementation and some limitations for a few of them, mostly on the Meta SDK. But the same parameter values gave visibly different output depending on the host, and every technique had to be retuned per platform.

\pointto{Right panel.} The dividing line is not how complex a technique is. It is whether the technique needs new parameter values or a new input signal. New values you can always find. A new signal you only get if the pipeline is open.

\remind{The last panel is the contribution. Everything before it is evidence.\\[0.4em] If asked which limitations: the Meta SDK gives one fragment hook, no second pass, no depth buffer and no normal buffer. The hierarchical technique had to read a camera distance derivative instead of a real depth texture, and Kuwahara lost its smoothing step.\\[0.4em] Do not re-justify the Avaturn choice, slide 11 already did.}

\slidehead{18}{RQ2 Answer}{0:40}

The second question asked which technique, or which combination, best balances visual abstraction against facial and expressive legibility.

\pointto{Left panel.} On the initial texture, none of them succeeded. The edge operators read the baked lighting as structure, and the features I actually needed stayed unmarked.

\pointto{Middle panel.} After the preprocessing, on the less noisy texture, some of them did succeed. Sobel now finds the eyebrows and the lip line at a moderate threshold. The Gaussian prefilter became unnecessary, because it gave no improvement over plain Sobel once the texture was clean. Five of the eight went forward as candidates.

\pointto{Right panel.} And the parameter adjustment study settled the rest. Kuwahara and the hierarchical layer were both excluded by the ranking, which leaves three: the outline, X-Toon and Sobel.

The answer to RQ2 is therefore a composite, not a technique. And not a technique alone either, because without the texture preprocessing none of the edge operators got there.

\remind{Say \emph{parameter adjustment study}, not \emph{user study}. The main user study measured trust on the finished composite, it did not select techniques.\\[0.4em] The five candidates were V1, V2.2, V4, V5 and V6. Only give that list if asked.}

\slidehead{19}{Main User Study}{1:00}

In the main user study we used three scenarios and three conditions: a real person, a realistic avatar, and the NPR abstracted version of that same avatar. The animation and the voice of C2 and C3 were taken from the real human recording in C1, so only the rendering changes between them. \pointto{The three faces.}

The scenario order was fixed, and the character order was rotated across three study versions.

\pointto{The procedure.} We first displayed the decision vignette, which explains the scenario and helps the participant put herself in that situation. Before the video we asked for her own choice. Then she watched the advisor, and afterwards we asked the questions that assess trust, and where she places that condition in the uncanny valley through eeriness and likeability. Three rounds, one per scenario, and then the overall assessment at the end.

\pointto{Left column.} The measures are all established instruments. Eeriness comes from Ho and MacDorman, likeability from the Godspeed questionnaire, and the three trust dimensions from Mayer, with the item wording adapted from Canales. The final forced choice is Canales' paradigm too. Alpha .87 to .96. Fifty-two people started, eleven failed the attention check.

\remind{Densest slide in the talk. Follow the diagram with your hand, do not read the panels.\\[0.4em] Say ``C2 versus C3 is the critical pair'' out loud. The whole results section rests on it.\\[0.4em] Do not say you took the \emph{design} from Canales. You took the trust item wording and the forced-choice measure. The three conditions, three scenarios and the Latin square are your own.\\[0.4em] Backups 49 to 53: the final stylisation question, scenarios, Latin square, full vignettes.}

\slidehead{20}{Four Hypotheses, and All Five Subscales}{0:40}

Before the analysis of the results, we had four hypotheses.

\pointto{Left panel.} H1 predicted that the stylised avatar would feel less eerie than the photorealistic one. H2 predicted that stylisation would cost no trust against that same avatar. H3 predicted that within a person, more eeriness lowers trust and more liking raises it. And H4 predicted no trust cost against the real human.

H2 and H4 predict a \emph{null}, and a non-significant p does not prove two things are the same. Both therefore carry an equivalence test, specified in advance.

\pointto{Right panel.} And this is every subscale at once. The next slides go through it one hypothesis at a time.

\remind{Give each hypothesis in one clause. Do not read the results off the figure here, the H slides do that.\\[0.4em] The sentence you are saving for later is that the real person separates from both avatars on all five, and the two avatars separate from each other on none of them. Say it on slide 26, not here.}

\vspace{0.6em}

%% ============================================================
\section*{\color{knblue}Part 4 — Results (slides 21 to 26, about 4 minutes)}

\slidehead{21}{H1 — Does NPR reduce eeriness?}{0:35}

My expectation was that the stylised avatar would be rated less eerie. It was not. Eeriness of the NPR avatar was slightly higher than the realistic one, and although that difference is not significant, \textbf{H1 is not supported}.

And this is not because people failed to notice the stylisation. The final stylisation question asked participants to rate how visually stylised each advisor looked, and it separated all three conditions. They saw the difference. They just did not find it any less eerie.

\remind{Say \emph{the final stylisation question}, not \emph{the manipulation check}.\\[0.4em] Backup slide 52 has the numbers if anyone wants them: omnibus $F(2,80)=15.45$, $p<.001$, and C2 versus C3 at $p=.004$.\\[0.4em] The effect between the two avatars is $g=0.07$, essentially zero. Give it only if asked how close they were.}

\slidehead{22}{H2 — Does NPR cost trust against the photorealistic avatar?}{0:45}

H2 predicted that stylisation would cost no trust against the photorealistic avatar.

The direct comparisons found nothing. Ability, benevolence and integrity are all non-significant, and all three effects are small. \textbf{So H2 is supported as stated}: no trust dimension came out lower for the stylised avatar.

But I also ran an equivalence test, because a non-significant result on its own does not prove two things are the same. That test does not clear it. The interval still contains both a small drop and a moderate gain, which means the two avatars are indistinguishable rather than demonstrably equivalent.

That is a sample size problem, not a surprise. At forty-one, the power for that test is about 0.2 even if the true difference were exactly zero.

\remind{If asked why you did not collect more: the deadline set the collection window and it fell in the semester break. Limitations slide, 30.\\[0.4em] Settling it properly needs roughly three times the participants.}

\slidehead{23}{H3 — Eeriness, likeability and trust}{0:40}

H3 predicted that trust correlates positively with likeability and negatively with eeriness.

In this analysis we did not compare between conditions. We asked instead: when the same person felt more eerie about one video than about their own average, did their trust drop.

And according to the results, it did. One point of eeriness costs about one point of trust, on all three trust dimensions, and likeability does the same in the opposite direction. \textbf{H3 is confirmed in both directions.}

However, when we check the correlation between those two, it is minus .72. They are highly correlated, which makes it hard to separate them as independent variables. When both go into the same model, likeability absorbs most of the eeriness effect. So I can say that appearance predicts trust. I cannot say which of the two is doing the work.

And this is the important one. The route from appearance to trust is open and wide. My composite never travelled it, because it never moved eeriness in the first place. That is why H1 and H2 are both null, and they fail for the same reason.

\slidehead{24}{H4 — Does NPR cost trust against the real speaker?}{0:40}

H4 predicted no trust cost against the real speaker.

There was one. \pointto{The three panels.} On all three trust dimensions the real person sits significantly above the stylised avatar, with large effects. You can read it off the brackets: three stars between the real speaker and the NPR avatar, and \emph{ns} between the two avatars.

\pointto{Second click, the equivalence plot appears.} And the equivalence test is nowhere close either. The shaded band is the bound I set in advance. The interval sits entirely outside it, to the right. The gap is 1.31 points, nearly three times that bound, so this is not a near miss. \textbf{H4 is not supported.}

Two explanations stay open and this design cannot separate them. Either this is what the uncanny valley predicts, or some participants recognised the speaker, because the speaker is me.

\remind{Build the slide in two clicks: the three panels first, then the equivalence plot. Explain each as it appears, and do not let the plot sit on screen unexplained while you are still on Step 1.\\[0.4em] Same two-beat structure works on H2, where the point is the opposite: the interval overlaps the band but still runs past it on one side.}

\slidehead{25}{The behavioural measure}{0:35}

Participants also had to choose one advisor rather than rate them. I included it because Canales found that what people \emph{say} and what people \emph{choose} can come apart, and in their data the stylised avatar won on choice.

In mine they did not come apart. The rankings reproduce the ratings. Real person wins, the two avatars tie.

\pointto{Right panel.} And before anyone asks whether the null is an artefact: I checked order, suspicion, the secondary measures and dispositional trust. None of them changes the picture.

\remind{Do not walk the four checks one by one unless asked. The panel is there for the examiner who wants it.}

\slidehead{26}{RQ3 answer}{0:40}

Back to the third question. How does stylisation affect perceived eeriness and trustworthiness, and how far does it reach.

Three things did not happen. Eeriness did not fall, it rose slightly. The two avatars were indistinguishable on all five subscales. And there was no behavioural advantage either.

One thing did. Stylisation carried no trust penalty I could detect.

\pointto{The box.} The contrast worth remembering: between the two render styles, a third of a point, inside the noise. Between the stylised avatar and the real human, 1.31, nearly three times the bound.

\remind{Say \emph{no penalty was detected}, not \emph{there is no penalty}. And \emph{indistinguishable}, not \emph{equivalent}.}

\slidehead{27}{Discussion}{0:03}

\remind{Section divider. Click straight through, do not read it out.}

\slidehead{28}{Why the two avatars scored the same}{0:55}

Two patterns need explaining. The real person sits far above both avatars everywhere. And the two avatars looked different to participants but scored identically on everything that mattered.

H3 already showed the mechanism was available. My composite changed what people saw without using it. Two mismatches inside the composite account for that.

First, material versus shape. My techniques change the surface, not the geometry. Zell and colleagues found material drives eeriness, which should have helped. But the same paper found eeriness goes \emph{up} when material and shape are stylised to different degrees, and that is exactly what C3 is: a cartoon surface on a face rebuilt from photographs. I may have removed one source of eeriness and added another.

Second, appearance versus performance. Both avatars use the same recorded human voice, gaze and lip movement, so C3 is a fully human performance coming out of a visibly drawn face. Cihodaru-\c{S}tefanache and Podina argue uncanniness comes from features perceived at inconsistent levels of humanness. A shading filter only changes appearance, so it could only widen that.

\remind{These two are grounded in published work. Do not rush them, this is where you sound like a researcher.}

\slidehead{29}{What this design could not show}{0:40}

Two reasons the effect could be real and still not appear here.

The comparison set. Every participant saw a real human in the same session. Once a real person anchors the top of the scale, the only question left is how far below. Both my avatars were obviously CG, so a rendering difference competes with a categorical one.

The setting. Flat screen, prerecorded, no interaction. Zibrek and colleagues found render style significantly moves presence in VR. What is too weak in a browser window may register once the avatar is in your room.

Neither is a defect in the analysis. Both say the same thing: this design tests stylisation under conditions that work against it.

\remind{Land on that last sentence. It turns a null result into a bounded one, and it sets up both limitations and future work.}

\slidehead{30}{Limitations}{0:40}

Six honest ones.

Session length. Both studies had to stay under fifteen minutes with every measure repeated three times, which capped what I could ask.

One configuration. One composite, one avatar, one platform. This describes \emph{this} configuration, not NPR as a whole. And the parameter study only tested the highest-impact groups, so the composite ranked highest among those tested rather than being the best available.

Sample size. Forty-one detects a difference, it does not prove equivalence. Counterbalancing. A full design needs six orders and twice the participants. I used three.

And one that matters more than the others. The three version groups differed significantly on eeriness and likeability, and eeriness is what H1 rests on. Version determines both who was in a group and the order they saw the conditions, and this design cannot separate the two. That ambiguity qualifies the H1 null.

\remind{Say the last one properly. Volunteering the limitation that cuts closest to your own null is what makes the rest of the talk credible.}

\slidehead{31}{Conclusion}{0:45}

To close. This thesis built NPR stylisation for three avatar platforms, eight techniques in sequence, each one answering a limitation of the one before it, narrowed to three by a visual assessment and the parameter adjustment study.

What the study showed is a split result. The composite did not reduce eeriness against the photorealistic baseline. But it produced no observed trust penalty either. The evidence does not indicate that stylisation reduced trust, and it is also not sufficient to establish that the two are equivalent.

What that means in practice is this. A developer choosing between the two styles need not assume an inherent disadvantage for the stylised one. And the gap that matters is not between the two render styles at all. It is the distance to the real human, and neither avatar condition closed any part of it.

Reaching the credibility of an actual human is a harder challenge than selecting a rendering style, and it is the problem that decides whether virtual advisors become genuinely trusted participants rather than a passing novelty.

\remind{Do not recap RQ1 to RQ3 here. They each already have an answer slide, and repeating them is the fastest way to flatten the ending.\\[0.4em] Land on the last sentence and stop. Do not add anything after it.}

\slidehead{32}{Future work}{0:35}

Five steps, in the order the thesis sets them out.

First, a larger replication. Repeat C2 versus C3 with about one hundred and twenty people, which settles the equivalence question either way.

Then the two explanations can each be tested directly. Stylise the geometry as well, so proportion shifts alongside the shading. And drop the real speaker, which removes the anchor and shows whether the two avatars separate with no human in view.

Fourth, sweep abstraction continuously instead of fixing it at one value, which replaces my single point with a curve and could reveal an optimum between full realism and full abstraction.

And fifth, restore what video removed: reciprocal behaviour and shared space in VR, across more than one avatar identity.

\remind{Give one clause each and move. Five items at four seconds apiece is the whole slide.\\[0.4em] The smaller fix on the bar, stating up front that the videos were not AI-generated, is worth naming only if you have time or if the generative-AI question came up earlier.}

\slidehead{33}{Thank you}{0:10}

That is the work. Eight techniques, three platforms, two studies, eighty-six participants in total. Thank you very much, and I am happy to take your questions.

\remind{Two slides sit behind this one and you do not present either: \textbf{29 References} and \textbf{30 Media credits}. Leave the thank-you slide up during questions. Jump to 29 only if someone asks which paper a claim came from.}

\clearpage

\clearpage

%% ============================================================
\section*{\color{knblue}Backup slides — not presented}

\remind{The backup section sits after the media credits and runs in talk order: the evidence timeline, then V1 to V6 across platforms, then the pre-study, then the study materials, then the extra statistics. Only one of them has a script, because only one of them you would ever \emph{present} rather than point at.}

\slidehead{B1}{Backup — How the evidence developed}{0:30}

In chronological order, that evidence falls into five stages, and the shape it makes is the point.

The first row is the valley being established, from Mori in 1970 through to Dill in 2012. By the end of that stretch we know it applies to CG characters, not only to robots.

Then 2012 is the turning point. McDonnell measures trustworthiness directly and finds it valley-shaped, and that is the first real evidence that abstraction is a way out rather than just a style choice.

2015 adds the condition. It is the material that drives eeriness, and mixing stylisation levels makes it worse.

And then look at the last two rows. Since 2021 the evidence splits. Canales and Dubosc point towards abstraction. Seymour, Gao and Tao point the other way, and Alipour's review of fifty-three studies says there is no consistent link at all.

That split is why this thesis exists.

\remind{This is the one backup worth walking through if a reviewer asks you to place the work in time. One sentence works if you are short: ``laid out in time, the evidence splits after 2021, and it has not converged since.''\\[0.4em] The other seventeen backups are reference material, not scripts. Know where they are: V1 to V6 on 38 to 45, the pre-study on 46 and 47, the study materials on 48 to 52, the extra statistics on 53 and 54.}

%% ============================================================
\section*{\color{knblue}Likely questions, and short answers}

\remind{Answer in two or three sentences, then stop. If you do not know, say what your data can and cannot show. That is a strong answer, not a weak one.}

\textbf{``Your main result is a null. Why should we believe it means anything?''}\\
Because I tested it two ways. The direct comparison found nothing, and I also pre-specified an equivalence test rather than treating a non-significant p-value as proof. H2 itself holds, but the equivalence test is inconclusive at this sample size, and I report that rather than claiming equivalence. And the final stylisation question shows participants clearly saw the difference, meaning this is not a failed stimulus.

\textbf{``Forty-one participants is small. Is the study underpowered?''}\\
For detecting a difference it met the power analysis. For proving equivalence it does not, and I report that. Equivalence needs roughly three times as many people. The collection window was fixed by the submission deadline and it fell in the semester break, when the mailing list channel was mostly empty.

\textbf{``You were the speaker in C1. Is that not a confound?''}\\
Yes, and I flag it as one. Some participants may have recognised me, and that gives C1 a basis for trust that no rendering can match. It inflates the C1 baseline, and the H4 gap may therefore be larger than it would be with a stranger. It does not affect the C2 versus C3 comparison, which is the critical one.

\textbf{``Why not just use one platform, if the study only used Avaturn?''}\\
Because RQ1 is a separate contribution. Building the same eight techniques three times is what exposed the difference between needing new parameter values and needing a new input signal. That is not visible from one platform. Avaturn was then selected for the study on three stated criteria, and it is the only one that meets all three.

\textbf{``Why did you exclude Kuwahara and the hierarchical layer?''}\\
Not by my own judgement. In the parameter study participants significantly preferred the version without them, in both cases. Since trustworthiness is the outcome I care about, that was the deciding evidence.

\textbf{``Is the composite optimal?''}\\
No, and I do not claim that. I tested the parameter groups with the largest expected visual impact and fixed the rest by inspection during development. It is a configuration judged suitable within the limits of the study, not a global optimum.

\textbf{``Could a stronger or different stylisation have worked?''}\\
Possibly. That is exactly why my main future work item is to sweep abstraction strength continuously instead of testing one setting. My data is a single point on a curve I never measured.

\textbf{``What is the practical takeaway for a developer?''}\\
Two things. Stylisation did not cost trust here, meaning a stylised avatar is not an obvious disadvantage. And if you are working on a closed platform, plan around the input signals it exposes, not around the techniques you would like to use.

\end{document}
